\section{Introduction}\label{rec:intro}

A recursively defined norm can be easy to evaluate at one node while
its estimates deteriorate as the tree grows. We study a recursion in
which two nonnegative fields encode a vector as their difference.
Each field at a node dominates the norm of its values at the children,
and the recursion chooses the least feasible pair. For Euclidean child
aggregation, one of the two inequalities is consequently tight at
every node. We show that this local equality controls midpoint
perturbations uniformly in the height and branching of the tree.

For an infinite-dimensional real Banach space $X$ with norm $N$, let
$\operatorname{cof}(X)$ denote its closed finite-codimensional linear
subspaces. The averaged asymptotic midpoint modulus is
\begin{equation}\label{intro:modulus}
 \widehat\delta_N(t)=\inf_{N(x)=1}\sup_{F\in\operatorname{cof}(X)}
 \inf_{\substack{y\in F\\N(y)\ge1}}
 \left(\frac{N(x+ty)+N(x-ty)}2-1\right),\qquad t>0.
\end{equation}
It measures the mean increase in the endpoint norms when a perturbation
avoids a suitable finite-dimensional part of the space. The usual
asymptotic midpoint modulus $m_N(t)$ uses the maximum of the two
endpoint norms, with unit directions, in the same infimum and supremum.
The norm is \emph{asymptotically midpoint uniformly convex}, abbreviated
AMUC, if $m_N(t)>0$ for every $t>0$. This condition was introduced by
Dilworth, Kutzarova, Randrianarivony, Revalski, and Zhivkov
\cite[Definition~2.2]{DKRRZ2016}.
The one-sided modulus is
\begin{equation}\label{intro:auc}
 \overline\delta_N(t)=\inf_{N(x)=1}\sup_{F\in\operatorname{cof}(X)}
 \inf_{\substack{y\in F\\N(y)=1}}\bigl(N(x+ty)-1\bigr).
\end{equation}
A norm is asymptotically uniformly convex, abbreviated AUC, if this
last quantity is positive for every $t>0$.

The averaged and maximum midpoint moduli have the same positivity
property, with
\begin{equation}\label{intro:comparisons}
 \tfrac12m_N(t)\le\widehat\delta_N(t)\le m_N(t).
\end{equation}
Indeed, average is at most maximum. For the converse, at a fixed unit
center intersect a witnessing subspace with the kernel of a norming
functional. Both endpoint norms are then at least one, and their
average excess is at least half the maximum excess. Taking the
infima and suprema proves \eqref{intro:comparisons}.
For a fixed direction, the mean of the two endpoint norms is an even
convex function of the radius. It is therefore nondecreasing on the
positive half-line. Thus the inner infimum in \eqref{intro:modulus}
can equally use $N(y)=1$, and $\widehat\delta_N$ is nondecreasing.
In particular, positivity for $0<t<1$ implies positivity at every radius.

The distinction from AUC renormability is already established.
Dilworth and coauthors separated the two norm properties on $\ell_2$
and posed the corresponding renorming question
\cite[Theorem~2.4 and Section~5]{DKRRZ2016}.
Baudier answered the general question negatively by using the
Kadets--Werner modification of the Bourgain--Rosenthal construction
\cite[Theorem~1 and Corollary~1]{Baudier2026}; the underlying construction
is due to Kadets and Werner \cite[Theorem~2.5]{KW2004}.
That example is nonreflexive \cite[Corollary~2.6]{KW2004}.
Our focus is the quantitative geometry of norms given by an explicit
tree recursion. The joining-root versions are reflexive and still
admit no equivalent AUC norm.

Tree norms have a substantial history in Banach-space theory.
James's construction uses sums over disjoint segments of a binary
tree \cite[p.~739]{James1974}; Girardi proved AUC for its canonical
predual and full dual \cite[Theorems~3 and~5]{Girardi2001}.
A closer comparison is the positive tree cone recalled by Banakh,
following Ghoussoub, Maurey, and Schachermayer
\cite[Section~2.I, pp.~28--29]{Banakh2000}: a nonnegative field at a
node dominates the Euclidean norm of its children. That construction
uses interpolation after taking a symmetric convex hull. Here the
vector is the difference of two least such fields, and the norm is an
explicit cost at the root. The defining norm and its completion are
proved directly below.

\paragraph{Results and proof structure.}
Section~\ref{rec:norms} defines three Euclidean constructions, proves
their completion and finite-coordinate projection properties, and
states Theorem~\ref{rec:main}. The first norm takes the sum of the
root fields on the infinite tree. The second joins finite-height components at
a zero coordinate and takes the common root field; it satisfies a
homogeneous cubic inequality. The third keeps a free coordinate at
that joining root and takes the sum. Both root-sum norms have averaged
modulus at least $t^3/128$ for $0<t<1$; the joined space is reflexive.
Distinguishing the root
normalizations is essential to the constants.

The shortest quantitative argument is in Section~\ref{rec:energy}.
Averaging the fields at $x+z$ and $x-z$ increases both fields of $x$
by the same nonnegative amount $d_s$ at node $s$. Writing $P_s,Q_s$ for
the fields of $x$, the expression
\[
 [d_s(2P_s+d_s)(2Q_s+d_s)]^{2/3}
\]
dominates the sum of its values at the children, by H\"older's
inequality and the tight child inequality. Iterating to the first
nodes outside a finite initial set containing the support of the center
controls the complete tail
by the defect at the root. No factor is multiplied once per level.
Section~\ref{rec:obstruction} supplies the bounded-path argument that
excludes AUC renormings.

Section~\ref{rec:martingale} gives a second way to propagate the same
root-sum defect. A finite stopping rule separates paths according to
the smaller field, a product of field ratios, and the local increase.
An absorbing martingale controls the small-product event. The proof
is uniform over the finite head and retains the order of its parameter
limits. Section~\ref{log:section} gives a different argument on the
joined tree: a logarithmic transform of two bounded supermartingales
bounds a product of ratios and proves a gap for infinite separated
midpoint families.
The stopping proof uses a predictable martingale transform, a classical
tool studied by Burkholder \cite{Burkholder1966}. In the logarithmic proof,
the selected supermartingale increments split into predictable drift and
centered martingale differences. The exact
finite square identities and stopping estimates used here are proved
in place. These arguments explain distinct mechanisms for uniformity, even
though the cubic bound already proves positive midpoint modulus.

Finally, Section~\ref{var:section} replaces Euclidean aggregation
inside a height-$h$ component by $\ell_{1+1/h}$ aggregation. This
exponent is constant inside each component. Two propagation arguments
control different energies: normalized iteration yields a sixth-power
estimate, while accumulated layer energy yields a cubic estimate with
coefficient $(6e)^{-3}$. Their complete proofs include the outer
$\ell_2$ sum and the precise finite-head and weakly null tail statements.
The cubic estimate also gives a cubic lower bound for the averaged
midpoint modulus of this variable-exponent space.
